\bookchapter{บทที่ 6: คุ้มครองชีวิตดิจิทัลของคุณ — ความปลอดภัยและความเป็นส่วนตัว}

\chapterintro{ความปลอดภัยของโทรศัพท์ = ความปลอดภัยของความเป็นส่วนตัว + ทรัพย์สิน ในยุคที่ข้อมูลคือทรัพย์สิน คุ้มครองโทรศัพท์ก็คือคุ้มครองตัวคุณเอง}

\booksection{หนึ่ง ล็อกหน้าจอ: แนวป้องกันทางกายภาพด่านแรก}

ล็อกหน้าจอคือแนวป้องกันด่านแรกของโทรศัพท์ และเป็นสิ่งที่คนมองข้ามง่ายที่สุด

\booksubsection{จัดอันดับระดับความปลอดภัย:}

\begin{enumerate}
  \item \textbf{รหัสผ่าน/รหัส PIN}: ปลอดภัยสูงสุด เป็นเกราะสุดท้าย (ใช้เมื่อลายนิ้ว/หน้าไม่ได้ผล)
  \item \textbf{ลายนิ้วมือ}: สมดุลที่ลงตัวระหว่างความปลอดภัยกับความสะดวก
  \item \textbf{การรู้จำใบหน้า}: สะดวกที่สุด แต่ความปลอดภัยขึ้นกับเทคโนโลยี (โครงสร้างแสง 3D ปลอดภัย ภาพ 2D หลอกได้ง่าย)
  \item \textbf{รูปแบบลากเส้นปลดล็อก}: สะดวกแต่มักทิ้งร่องรอยจนถูกเดาได้
\end{enumerate}

\booksubsection{ตั้งล็อกหน้าจอให้แน่นหนา:}

\begin{itemize}
  \item ตั้งรหัสผ่านที่ซับซ้อนให้ได้เลย (อย่าใช้วันเกิด หรือ 123456)
  \item เปิดการยืนยันด้วยสิ่งมีชีวิต (ไบโอเมตริกซ์)
  \item ตั้งเวลาล็อกหน้าจอเป็น "ทันที" หรือสั้นที่สุด
\end{itemize}

\booksection{สอง “ค้นหาโทรศัพท์ของฉัน”: อุปกรณ์กันขโมยขั้นสุดท้าย}

\textbf{ต้องเปิดไว้!}มันทำอะไรได้:

\begin{itemize}
  \item \textbf{ระบุตำแหน่ง}: โทรศัพท์หายไปไหน หาบนแผนที่ได้
  \item \textbf{เล่นเสียง}: ทำให้โทรศัพท์ดังขึ้น เหมาะกับการหาในบ้าน
  \item \textbf{โหมดสูญหาย}: ล็อกเครื่องจากระยะไกลพร้อมโชว์เบอร์ติดต่อ
  \item \textbf{ลบข้อมูลจากระยะไกล}: ท่าสุดท้าย ล้างข้อมูลทั้งหมด (แต่ล้างแล้วระบุตำแหน่งต่อไม่ได้)
\end{itemize}

ไปที่ [ตั้งค่า] ฝั่งแอปเปิลหา “ค้นหา iPhone” ฝั่งแอนดรอยด์หา “ค้นหาอุปกรณ์ของฉัน” หรือบริการคล้ายกันของแบรนด์ ให้แน่ใจว่ามันเปิดอยู่ (สถานะสีเขียว)

\booksection{สาม การอัปเดตระบบ: เย็บแพตช์ อุดช่องโหว่}

เป้าหมายแรกของการอัปเดตระบบคือแก้ช่องโหว่ความปลอดภัย เหมือนอุดรูบนกำแพงบ้าน เห็นแจ้งเตือนอัปเดตระบบเมื่อไร พยายามอัปเดตทันทีเมื่อไร

แน่นอนว่าการอัปเดตเวอร์ชันใหญ่ (เช่น iOS 17 ไป iOS 18) อาจมีปัญหาความเข้ากันได้ รอดูอีกสองสามวันค่อยอัปก็ได้

\booksection{สี่ รู้เท่าทันและป้องกันมิจฉาชีพทางโทรศัพท์}

ในยุคที่ AI ก้าวหน้ารวดเร็ว ความเชื่อเดิมว่า “ตาเห็นคือจริง” ใช้ไม่ได้อีกแล้ว มิจฉาชีพกำลังใช้เทคโนโลยี Deepfake และ AI ยกระดับจากการ “โยนแหกับดักกว้างๆ” มาเป็นการ “ล่าเป้าแบบเจาะจง” อย่างครบวงจร

\booksubsection{กลโกงยุค AI ที่พบบ่อยที่สุด}

\textbf{(1) กลโกงสลับหน้าและเลียนแบบเสียงด้วย AI}\par

พวกมิจฉาชีพเก็บรูปความละเอียดสูงที่เป้าหมายโพสต์ในโซเชียล หรือใช้เสียงตัวอย่างแค่ไม่กี่วินาที แล้วใช้เทคโนโลยี Deepfake โคลนเป็น “วิดีโอ/เสียง” ของเพื่อนญาติที่เหมือนจริงจนแยกไม่ออก บทที่ใช้บ่อยเช่น\textbf{รถชนฉับพลัน ต้องยืมเงินด่วน ถูกลักหัวปิด}ฯลฯ กลโกงแบบนี้อาศัยความผูกพันทางครอบครัวและมิตรภาพ\textbf{หลอกลวงสูงมาก}.

\textbf{(2) แอบอ้างหน่วยงาน / ฝ่ายบริการลูกค้า}\par

เทคโนโลยี AI เลียนแบบน้ำเสียงและโทนของเจ้าหน้าที่ทางการ ตำรวจ อัยการ ศาล ได้แม่นยำ แม้แต่ตัดต่อภาพหน้าจอโอนเงินปลอม เอกสารราชการ หมายเรียก ฯลฯ ที่สมจริงสูง ผ่านการ\textbf{สร้างความตื่นตระหนกหรือกดดันด้วยอำนาจ}จนหลอกให้เหยื่อโอนเงินหรือเปิดแชร์หน้าจอ เพื่อขโมยข้อมูลบัญชี

\textbf{(3) กลโกงลงทุนด้วย AI}\par

ใช้อัลกอริทึม AI ปลอมภาพหน้าจอผลกำไร “กำไรไม่มีขาดทุน” หน้าโปรแกรมเลือกหุ้นเชิงปริมาณปลอม หรือรายงานวิเคราะห์การลงทุน “ระดับกูรู” เอาผลตอบแทนสูงเป็นเหยื่อล่อ ขายบริการลงทุนราคาแพง หรือหลอกให้เหยื่อโอนเงินเข้าบัญชีเทรดที่พวกมันควบคุม\textbf{ท้ายที่สุดก็หอบเงินหนีไป}.

\booksubsection{ทางแก้ — แนวคิดป้องกันแบบ “กระจกเงาส่องปีศาจ”}

AI โคลนเสียงและใบหน้าได้ แต่ก็ไม่ใช่ไม่มีจุดอ่อน มีรอยต่อให้ใช้ประโยชน์อยู่:

\begin{itemize}
  \item \textbf{วิธียืนยันด้วยการเคลื่อนไหว}: ตอนวิดีโอคอล ขอให้อีกฝ่ายทำท่าทางต่อเนื่องที่นัดไว้ เช่น "เอามือโบกผ่านหน้าเร็วๆ" "เอียงหน้าให้เห็นใบหู" หรือ "กระพริบตาเร็วๆ 5 ครั้ง" เทคโนโลยีสลับหน้า AI ปัจจุบันเจอฉากบังหน้าซับซ้อน การเคลื่อนไหวเร็ว และแสงเงาที่เปลี่ยนแล้ว\textbf{มักภาพกระตุก ใบหน้าบิดเบี้ยว หรือขอบทะลุเฟรม}ง่ายมาก
  \item \textbf{วิธียืนยันด้วยรหัสลับ}: สร้างและใช้คำถามยืนยันส่วนตัวที่นัดกันไว้ เช่น "วันที่เราไปกินข้าวด้วยกันครั้งล่าสุดนายสั่งอะไร" "ของขวัญชิ้นแรกที่ฉันให้นายคืออะไร" รายละเอียดที่รู้กันเฉพาะสองคน ถ้าอีกฝ่ายลังเลชัดๆ หรือตอบไม่ตรงคำถาม ให้\textbf{วางสายทันทีและเปิดแผนรับมือความปลอดภัย}.
  \item \textbf{วิธียืนยันซ้ำสองชั้น}: ยึดหลัก "สี่ไม่หนึ่งตรวจ" ที่หน่วยงานความมั่นคงไซเบอร์ประกาศอย่างเคร่งครัด —\textbf{ไม่เชื่อข้อมูลแปลกปลอมง่ายๆ ไม่โอนเงินตามใจลอย ไม่เปิดเผยข้อมูลอ่อนไหวส่วนตัว ไม่แชร์เนื้อหาที่น่าสงสัย}และ\textbf{เมื่อไหร่ที่เกี่ยวกับเงินทอง ต้องโทรกลับด้วยเบอร์เดิมที่รู้จัก หรือยืนยันตัวต่อหน้า}. ตอนโอนเงิน แนะนำเลือกโหมด "เงินเข้าใน 24 ชั่วโมง" เผื่อเวลาให้การรู้เท่าทันกลโกงและการเรียกคืนเงิน
\end{itemize}

\booksubsection{กลโกงคลาสสิก (ยังต้องระวัง):}

\begin{itemize}
  \item \textbf{รางวัลล็อตเตอรี่/ส่งต่อคอมมิชชั่น}: เอาประโยชน์ล่อ ให้หวานนิดเดียวแล้วเอาเงินก้อนใหญ่
  \item \textbf{กลโกงหวาน}: กับดักหมู ปั้นความสัมพันธ์ก่อนแล้วค่อยโกงเงิน
  \item \textbf{แอบอ้างเป็นคนรู้จัก}: ขโมยบัญชีแล้วยืมเงิน (ยุค AI แยกยากขึ้นอีก)
\end{itemize}

\booksubsection{ระบบป้องกันในตัวโทรศัพท์:}

\begin{itemize}
  \item บล็อกสายรบกวน
  \item ป้องกันรหัส OTP ถูกขโมย
  \item ตรวจสอบความปลอดภัยของลิงก์ใน SMS
\end{itemize}

\booksubsection{คาถากันโกง:}

\textbf{ไม่เชื่อง่าย ไม่เปิดเผย ไม่กดลิงก์ ไม่โอนเงิน มีข้อสงสัยตรวจสอบก่อน!}เจอ “เรื่องด่วน” ใจเย็น 5 นาทีก่อน กลโกงแทบทุกแบบจะแฉตรงนี้แหละ

\booksection{ห้า รับมือ “กล่องเปิดข้อมูล” และการโจมตีทางสังคม: ตัดสายการรั่วไหลของความเป็นส่วนตัว}

หลักการแก่นของ “การเปิดกล่อง” (การโจมตีทางวิศวกรรมสังคม) คือ\textbf{เก็บรวบรวม วิเคราะห์ และประกอบร่างเศษข้อมูลดิจิทัลของบุคคลอย่างเป็นระบบ จนได้ภาพร่างของบุคคลนั้นมาทั้งชุด}หัวใจของการป้องกันคือวางเส้นแดงให้ความเป็นส่วนตัวชัดเจน และสร้างระบบป้องกันหลายชั้น

\booksubsection{การถอดข้อมูลออกจากรูปภาพ}

วิธีล้าง EXIF พื้นฐานดูบทที่ 2 “การจัดการรูปภาพและระบบไฟล์” ก่อนแชร์รูปลงโซเชียล ยังต้อง\textbf{ลบข้อมูลเมตา (EXIF) ของรูปให้หมดจริงๆ}ด้วย ทั้งเวลาถ่าย รุ่นอุปกรณ์ พิกัด GPS ฯลฯ ข้อมูลอ่อนไหว และ\textbf{บ้านเลขที่ ทะเบียนรถ บาร์โค้ดตั๋ว เลขบัตรประชาชน}ข้อมูลสำคัญในรูปต้องระบายทับหรือเบลอให้หมด กันเผลอเปิดเผยเส้นทางการเดินทางส่วนตัว

\begin{itemize}
  \item เข้าไปที่ iOS: [ตั้งค่า] > [ความเป็นส่วนตัวและความปลอดภัย] > [กล้อง] ปิด "ตำแหน่งที่แม่นยำ" ได้ ตอนแชร์เลือก "เฉพาะรูปที่เลือก" จะลบข้อมูลเมตาออกให้
  \item แอนดรอยด์: แอปอัลบั้มบางตัว (เช่น Google รูปภาพ) จะถามตอนแชร์ว่าจะเก็บข้อมูลตำแหน่งไว้ไหม
  \item เครื่องมือทั่วไป: ExifTool, Metapho ฯลฯ ดูและแก้ไข EXIF ได้
\end{itemize}

\booksubsection{การยุบสิทธิ์ของแอป}

วิธีจัดการสิทธิ์พื้นฐานดูบทที่ 4 “การจัดการสิทธิ์” ในบริบท “กันกล่องเปิดข้อมูล” ยังต้องไล่ตรวจและปิด\textbf{ไมโครโฟน สมุดโทรศัพท์ ข้อมูลตำแหน่ง}สิทธิ์ที่ไม่จำเป็นของแอปช้อปปิ้งและเครื่องมือเป็นระยะ กันแอป “แอบฟัง” บทสนทนาประจำวันหรือเก็บเครือข่ายความสัมพันธ์ส่วนตัวโดยผิดกฎหมาย แอปที่ทิ้งไว้นานแต่ยังถือสิทธิ์สูงควรถอนการติดตั้งทันที ลดจุดเสี่ยงการรั่วไหลตั้งแต่ต้นทาง

\begin{itemize}
  \item เห็นคำขอ "สิทธิ์สมุดโทรศัพท์" ให้ถามตัวเองก่อน: แอปนี้จำเป็นต้องรู้คนในสมุดโทรศัพท์ฉันทั้งเล่มเลยหรือ?
  \item แอปไหนสามเดือนไม่ได้ใช้ ถอนการติดตั้งเลย สิทธิ์ต้องเข้มขึ้นเรื่อยๆ ไม่ใช่หลวมลงเรื่อยๆ
\end{itemize}

\booksubsection{ตั้งกันสำหรับอุปกรณ์อัจฉริยะ}

กล้องวงจรปิดอัจฉริยะ ล็อกประตูอัจฉริยะ ลำโพงอัจฉริยะ ฯลฯ อุปกรณ์ IoT ตอนเริ่มใช้งานต้อง\textbf{เปลี่ยนรหัสผ่านเริ่มต้นทันที}ตั้งรหัสผ่านแยกที่แข็งแรง มีตัวพิมพ์ใหญ่เล็ก ตัวเลข และสัญลักษณ์พิเศษ ฟังก์ชันเข้าถึงจากระยะไกลที่ไม่ค่อยใช้ควรปิดไว้ ตรวจอัปเดตเฟิร์มแวร์เป็นระยะและแก้ช่องโหว่ให้ทัน\textbf{กันแฮกเกอร์บุกเข้ามาลอบดูพื้นที่ส่วนตัวของบ้าน}.

\begin{tcbwarn}{เคสจริงเตือนใจ}
ปี 2022 ตำรวจที่ไหนสักแห่งแจ้งว่า แฮกเกอร์เจาะรหัสผ่านเริ่มต้นของกล้องวงจรปิดบ้านแบรนด์หนึ่ง (หลายคนไม่เคยเปลี่ยน) ขโมยภาพห้องนอนของบ้านเกินพันหลังมาขายบนออนไลน์\textbf{ไม่เปลี่ยนรหัสเริ่มต้น = เปิดประตูบ้านกว้าง}.

\end{tcbwarn}

\booksection{หก คุ้มครองความเป็นส่วนตัว: ข้อมูลของคุณคือทรัพย์สินของคุณ}

รูปภาพ ประวัติแชท เส้นทางการเดินทางในโทรศัพท์ แท้จริงคือ “ทรัพย์สินดิจิทัล” ของคุณ พอรั่วไหลเบาก็โดนโฆษณารบกวน หนักก็โดนไล่ตามตัว กลโกง หรือแม้แต่แบล็กเมล์

\booksubsection{หลักการ “ลดให้น้อยที่สุด” ของสิทธิ์แอป:}

\begin{itemize}
  \item \textbf{ตำแหน่ง}: แอปเรียกรถ/แผนที่ใช้ "อนุญาตตลอด" แอปอื่นใช้ "อนุญาตเฉพาะขณะใช้งาน"
  \item \textbf{สมุดโทรศัพท์}: ให้เฉพาะเมื่อจำเป็น (เช่นเพิ่มเพื่อนใน WeChat) ไม่งั้นห้ามหมด
  \item \textbf{อัลบั้ม/รูปทั้งหมด}: iOS 14+ ขึ้นไปอนุญาต "รูปที่เลือก" ได้ แอนดรอยด์ 13+ ก็รองรับการอนุญาตแบบละเอียดเช่นกัน
  \item \textbf{ไมโครโฟน/กล้อง}: ปิดไว้เป็นปกติ เอาออกเมื่อต้องใช้ สงสัยว่าถูกแอบอัด? ปัดศูนย์ควบคุมออกมาดูว่ามี "จุดเขียว/จุดส้ม" ไหม
  \item \textbf{การติดตามข้ามแอปในพื้นหลัง}: iOS 14.5+ บังคับว่าแอปต้องถามก่อนใช้ IDFA แอนดรอยด์ปิด "รหัสโฆษณา" ได้
\end{itemize}

\booksubsection{การจัดการความเป็นส่วนตัวลงมือทำจริง:}

\begin{itemize}
  \item \textbf{รายงานความเป็นส่วนตัว (iOS)}: [ตั้งค่า] > [ความเป็นส่วนตัวและความปลอดภัย] > [รายงานความเป็นส่วนตัวของแอป] ดูว่าแอปไหนแอบเข้าถึงตำแหน่ง กล้อง ไมโครโฟนของคุณ
  \item \textbf{บันทึกการใช้สิทธิ์ (แอนดรอยด์)}: [ตั้งค่า] > [ความเป็นส่วนตัว] > [บันทึกการใช้สิทธิ์] ดูว่าใครเรียกใช้สิทธิ์อ่อนไหวเมื่อไร
  \item \textbf{จำกัดการติดตามเพื่อโฆษณา}: ปิด "รหัสโฆษณา" หรือ "จำกัดการติดตามโฆษณา" ในการตั้งค่าระบบ ช่วยลด "โดนโฆษณาเจาะจงตัว" ได้
\end{itemize}

\booksection{เจ็ด รายการ “ปฐมพยาบาล” หลังโทรศัพท์หาย/ถูกขโมย}

โทรศัพท์หายอย่าตกใจ ไล่จัดการตามลำดับนี้:

\begin{enumerate}
  \item \textbf{ระงับซิมทันที}: ติดต่อผู้ให้บริการขอระงับบัตร กันการเอาไปรับรหัส OTP
  \item \textbf{ล็อกจากระยะไกล}: เปิด "โหมดสูญหาย" ผ่านบริการคลาวด์ (iCloud / คลาวด์ของแบรนด์)
  \item \textbf{เปลี่ยนรหัสผ่านสำคัญ}: บัญชีสำคัญทั้งหมดที่เคยล็อกอินบนเครื่อง (ชำระเงิน ธนาคาร อีเมล) เปลี่ยนรหัสทันที
  \item \textbf{แจ้งความ}: ไปแจ้งความที่สถานีตำรวจ จดหมายเลข IMEI ไว้ เวลาเจอเครื่องก็พิสูจน์ตัวตนได้สะดวก
  \item \textbf{แจ้งญาติมิตร}: กันมิจฉาชีพยืมชื่อคุณไปยืมเงิน
  \item \textbf{(ท้ายสุด) ลบข้อมูลจากระยะไกล}: ทำเมื่อยืนยันแล้วว่าหาคืนไม่ได้ หลังล้างแล้วฟังก์ชันระบุตำแหน่งจะใช้ไม่ได้อีก
\end{enumerate}

\booksection{แปด ความปลอดภัยซิมกับ eSIM}

\begin{itemize}
  \item \textbf{ตั้ง PIN ซิม}: เปิดแล้วทุกครั้งที่รีบูตหรือเปลี่ยนซิมต้องกรอก PIN กันการ "ดึงซิมไปใส่เครื่องอื่น" ข้ามล็อกหน้าจอ
  \item \textbf{ข้อดีของ eSIM}: ไม่มีบัตรกายภาพให้ดึงออก ปลอดภัยกว่า พร้อมสลับหลายเบอร์ได้ สะดวกมากตอนเดินทาง
  \item \textbf{เสี่ยงการเปลี่ยนซิมซ้ำ}: ระวังกล "เปลี่ยนซิมซ้ำของผู้ให้บริการ" — มิจฉาชีพปลอมตัวเป็นคุณไปทำซิมใหม่ แล้วยึดเบอร์โทรศัพท์คุณ การตั้งให้ "ทำซิมใหม่ต้องเจ้าของเบอร์แบกบัตรประชาชนไปที่สาขาเอง" คือแนวป้องกันที่ได้ผลที่สุด
\end{itemize}

\begin{tcbwarn}{สิ่งที่ห้ามทำเด็ดขาด}
\begin{itemize}
  \item ใช้รหัสผ่านอ่อนแบบวันเกิด / 123456 / qwerty
  \item ส่งรหัส OTP ให้ใครก็แล้วกัน (แม้แต่ "เจ้าหน้าที่")
  \item กดลิงก์แปลกปลอมใน SMS โดยเฉพาะเบอร์ต่างประเทศที่ขึ้นต้น 00
  \item ล็อกอินธนาคารหรือชำระเงินบนไวไฟสาธารณะ
  \item เปิด "แชร์หน้าจอ" หรือ "ควบคุมระยะไกล" ให้คนแปลกหน้า
\end{itemize}

\end{tcbwarn}

\begin{tcbpractice}{ลองเปิดใช้และจำลอง "ล็อกจากระยะไกล"}
\begin{enumerate}
  \item มั่นใจว่าฟังก์ชัน "ค้นหาโทรศัพท์ของฉัน" เปิดอยู่
  \item บนเบราว์เซอร์คอมพิวเตอร์หรือโทรศัพท์อีกเครื่อง ล็อกอินบัญชีบริการคลาวด์ที่เกี่ยวข้อง (iCloud.com หรือเว็บคลาวด์ของแบรนด์)
  \item หาฟังก์ชัน "ค้นหาอุปกรณ์" ลองสั่ง "เล่นเสียง" หรือจำลองคำสั่ง "ล็อก" ใส่โทรศัพท์ของตัวเองหนึ่งครั้ง เพื่อคุ้นเคยกับกระบวนการทั้งหมด
\end{enumerate}

\end{tcbpractice}
